\IfFileExists{stacks-project.cls}{%
\documentclass{stacks-project}
}{%
\documentclass{amsart}
}

% For dealing with references we use the comment environment
\usepackage{verbatim}
\newenvironment{reference}{\comment}{\endcomment}
%\newenvironment{reference}{}{}
\newenvironment{slogan}{\comment}{\endcomment}
\newenvironment{history}{\comment}{\endcomment}

% For commutative diagrams we use Xy-pic
\usepackage[all]{xy}

% We use 2cell for 2-commutative diagrams.
\xyoption{2cell}
\UseAllTwocells

% We use multicol for the list of chapters between chapters
\usepackage{multicol}

% This is generally recommended for better output
\usepackage{lmodern}
\usepackage[T1]{fontenc}

% For cross-file-references
\usepackage{xr-hyper}

% Package for hypertext links:
\usepackage{hyperref}

% For any local file, say "hello.tex" you want to link to please
% use \externaldocument[hello-]{hello}
\externaldocument[introduction-]{introduction}
\externaldocument[conventions-]{conventions}
\externaldocument[sets-]{sets}
\externaldocument[categories-]{categories}
\externaldocument[topology-]{topology}
\externaldocument[sheaves-]{sheaves}
\externaldocument[sites-]{sites}
\externaldocument[stacks-]{stacks}
\externaldocument[fields-]{fields}
\externaldocument[algebra-]{algebra}
\externaldocument[brauer-]{brauer}
\externaldocument[homology-]{homology}
\externaldocument[derived-]{derived}
\externaldocument[simplicial-]{simplicial}
\externaldocument[more-algebra-]{more-algebra}
\externaldocument[smoothing-]{smoothing}
\externaldocument[modules-]{modules}
\externaldocument[sites-modules-]{sites-modules}
\externaldocument[injectives-]{injectives}
\externaldocument[cohomology-]{cohomology}
\externaldocument[sites-cohomology-]{sites-cohomology}
\externaldocument[dga-]{dga}
\externaldocument[dpa-]{dpa}
\externaldocument[sdga-]{sdga}
\externaldocument[hypercovering-]{hypercovering}
\externaldocument[schemes-]{schemes}
\externaldocument[constructions-]{constructions}
\externaldocument[properties-]{properties}
\externaldocument[morphisms-]{morphisms}
\externaldocument[coherent-]{coherent}
\externaldocument[divisors-]{divisors}
\externaldocument[limits-]{limits}
\externaldocument[varieties-]{varieties}
\externaldocument[topologies-]{topologies}
\externaldocument[descent-]{descent}
\externaldocument[perfect-]{perfect}
\externaldocument[more-morphisms-]{more-morphisms}
\externaldocument[flat-]{flat}
\externaldocument[groupoids-]{groupoids}
\externaldocument[more-groupoids-]{more-groupoids}
\externaldocument[etale-]{etale}
\externaldocument[chow-]{chow}
\externaldocument[intersection-]{intersection}
\externaldocument[pic-]{pic}
\externaldocument[weil-]{weil}
\externaldocument[adequate-]{adequate}
\externaldocument[dualizing-]{dualizing}
\externaldocument[duality-]{duality}
\externaldocument[discriminant-]{discriminant}
\externaldocument[derham-]{derham}
\externaldocument[local-cohomology-]{local-cohomology}
\externaldocument[algebraization-]{algebraization}
\externaldocument[curves-]{curves}
\externaldocument[resolve-]{resolve}
\externaldocument[models-]{models}
\externaldocument[functors-]{functors}
\externaldocument[equiv-]{equiv}
\externaldocument[pione-]{pione}
\externaldocument[etale-cohomology-]{etale-cohomology}
\externaldocument[proetale-]{proetale}
\externaldocument[relative-cycles-]{relative-cycles}
\externaldocument[more-etale-]{more-etale}
\externaldocument[trace-]{trace}
\externaldocument[crystalline-]{crystalline}
\externaldocument[spaces-]{spaces}
\externaldocument[spaces-properties-]{spaces-properties}
\externaldocument[spaces-morphisms-]{spaces-morphisms}
\externaldocument[decent-spaces-]{decent-spaces}
\externaldocument[spaces-cohomology-]{spaces-cohomology}
\externaldocument[spaces-limits-]{spaces-limits}
\externaldocument[spaces-divisors-]{spaces-divisors}
\externaldocument[spaces-over-fields-]{spaces-over-fields}
\externaldocument[spaces-topologies-]{spaces-topologies}
\externaldocument[spaces-descent-]{spaces-descent}
\externaldocument[spaces-perfect-]{spaces-perfect}
\externaldocument[spaces-more-morphisms-]{spaces-more-morphisms}
\externaldocument[spaces-flat-]{spaces-flat}
\externaldocument[spaces-groupoids-]{spaces-groupoids}
\externaldocument[spaces-more-groupoids-]{spaces-more-groupoids}
\externaldocument[bootstrap-]{bootstrap}
\externaldocument[spaces-pushouts-]{spaces-pushouts}
\externaldocument[spaces-chow-]{spaces-chow}
\externaldocument[groupoids-quotients-]{groupoids-quotients}
\externaldocument[spaces-more-cohomology-]{spaces-more-cohomology}
\externaldocument[spaces-simplicial-]{spaces-simplicial}
\externaldocument[spaces-duality-]{spaces-duality}
\externaldocument[formal-spaces-]{formal-spaces}
\externaldocument[restricted-]{restricted}
\externaldocument[spaces-resolve-]{spaces-resolve}
\externaldocument[formal-defos-]{formal-defos}
\externaldocument[defos-]{defos}
\externaldocument[cotangent-]{cotangent}
\externaldocument[examples-defos-]{examples-defos}
\externaldocument[algebraic-]{algebraic}
\externaldocument[examples-stacks-]{examples-stacks}
\externaldocument[stacks-sheaves-]{stacks-sheaves}
\externaldocument[criteria-]{criteria}
\externaldocument[artin-]{artin}
\externaldocument[quot-]{quot}
\externaldocument[stacks-properties-]{stacks-properties}
\externaldocument[stacks-morphisms-]{stacks-morphisms}
\externaldocument[stacks-limits-]{stacks-limits}
\externaldocument[stacks-cohomology-]{stacks-cohomology}
\externaldocument[stacks-perfect-]{stacks-perfect}
\externaldocument[stacks-introduction-]{stacks-introduction}
\externaldocument[stacks-more-morphisms-]{stacks-more-morphisms}
\externaldocument[stacks-geometry-]{stacks-geometry}
\externaldocument[moduli-]{moduli}
\externaldocument[moduli-curves-]{moduli-curves}
\externaldocument[examples-]{examples}
\externaldocument[exercises-]{exercises}
\externaldocument[guide-]{guide}
\externaldocument[desirables-]{desirables}
\externaldocument[coding-]{coding}
\externaldocument[obsolete-]{obsolete}
\externaldocument[fdl-]{fdl}
\externaldocument[index-]{index}

% Theorem environments.
%
\theoremstyle{plain}
\newtheorem{theorem}[subsection]{Theorem}
\newtheorem{proposition}[subsection]{Proposition}
\newtheorem{lemma}[subsection]{Lemma}

\theoremstyle{definition}
\newtheorem{definition}[subsection]{Definition}
\newtheorem{example}[subsection]{Example}
\newtheorem{exercise}[subsection]{Exercise}
\newtheorem{situation}[subsection]{Situation}

\theoremstyle{remark}
\newtheorem{remark}[subsection]{Remark}
\newtheorem{remarks}[subsection]{Remarks}

\numberwithin{equation}{subsection}

% Macros
%
\def\lim{\mathop{\mathrm{lim}}\nolimits}
\def\colim{\mathop{\mathrm{colim}}\nolimits}
\def\Spec{\mathop{\mathrm{Spec}}}
\def\Hom{\mathop{\mathrm{Hom}}\nolimits}
\def\Ext{\mathop{\mathrm{Ext}}\nolimits}
\def\SheafHom{\mathop{\mathcal{H}\!\mathit{om}}\nolimits}
\def\SheafExt{\mathop{\mathcal{E}\!\mathit{xt}}\nolimits}
\def\Sch{\mathit{Sch}}
\def\Mor{\mathop{\mathrm{Mor}}\nolimits}
\def\Ob{\mathop{\mathrm{Ob}}\nolimits}
\def\Sh{\mathop{\mathit{Sh}}\nolimits}
\def\NL{\mathop{N\!L}\nolimits}
\def\CH{\mathop{\mathrm{CH}}\nolimits}
\def\proetale{{pro\text{-}\acute{e}tale}}
\def\etale{{\acute{e}tale}}
\def\QCoh{\mathit{QCoh}}
\def\Ker{\mathop{\mathrm{Ker}}}
\def\Im{\mathop{\mathrm{Im}}}
\def\Coker{\mathop{\mathrm{Coker}}}
\def\Coim{\mathop{\mathrm{Coim}}}

% Boxtimes
%
\DeclareMathSymbol{\boxtimes}{\mathbin}{AMSa}{"02}

%
% Macros for moduli stacks/spaces
%
\def\QCohstack{\mathcal{QC}\!\mathit{oh}}
\def\Cohstack{\mathcal{C}\!\mathit{oh}}
\def\Spacesstack{\mathcal{S}\!\mathit{paces}}
\def\Quotfunctor{\mathrm{Quot}}
\def\Hilbfunctor{\mathrm{Hilb}}
\def\Curvesstack{\mathcal{C}\!\mathit{urves}}
\def\Polarizedstack{\mathcal{P}\!\mathit{olarized}}
\def\Complexesstack{\mathcal{C}\!\mathit{omplexes}}
% \Pic is the operator that assigns to X its picard group, usage \Pic(X)
% \Picardstack_{X/B} denotes the Picard stack of X over B
% \Picardfunctor_{X/B} denotes the Picard functor of X over B
\def\Pic{\mathop{\mathrm{Pic}}\nolimits}
\def\Picardstack{\mathcal{P}\!\mathit{ic}}
\def\Picardfunctor{\mathrm{Pic}}
\def\Deformationcategory{\mathcal{D}\!\mathit{ef}}
\usepackage{cleveref}
\crefname{originalreference}{original reference}{original references}
\Crefname{originalreference}{Original reference}{Original references}
\makeatletter
\newlabel{algebra-lemma-completion-flat}{{00MB}{source}{}{sourcepointer-1}{}}
\newlabel{algebra-lemma-completion-flat@cref}{{[originalreference][0][]00MB}{[1][0][]source}{}{}{}}
\makeatother
\crefname{originalreference}{original reference}{original references}
\Crefname{originalreference}{Original reference}{Original references}
\makeatletter
\newlabel{algebra-lemma-completion-tensor}{{00MA}{source}{}{sourcepointer-2}{}}
\newlabel{algebra-lemma-completion-tensor@cref}{{[originalreference][0][]00MA}{[1][0][]source}{}{}{}}
\makeatother
\crefname{originalreference}{original reference}{original references}
\Crefname{originalreference}{Original reference}{Original references}
\makeatletter
\newlabel{algebra-lemma-flat-increases-depth}{{00LM}{source}{}{sourcepointer-3}{}}
\newlabel{algebra-lemma-flat-increases-depth@cref}{{[originalreference][0][]00LM}{[1][0][]source}{}{}{}}
\makeatother
\crefname{originalreference}{original reference}{original references}
\Crefname{originalreference}{Original reference}{Original references}
\makeatletter
\newlabel{algebra-lemma-grothendieck-regular-sequence}{{00MG}{source}{}{sourcepointer-4}{}}
\newlabel{algebra-lemma-grothendieck-regular-sequence@cref}{{[originalreference][0][]00MG}{[1][0][]source}{}{}{}}
\makeatother
\crefname{originalreference}{original reference}{original references}
\Crefname{originalreference}{Original reference}{Original references}
\makeatletter
\newlabel{algebra-lemma-regular-quotient-regular}{{00NR}{source}{}{sourcepointer-5}{}}
\newlabel{algebra-lemma-regular-quotient-regular@cref}{{[originalreference][0][]00NR}{[1][0][]source}{}{}{}}
\makeatother
\crefname{originalreference}{original reference}{original references}
\Crefname{originalreference}{Original reference}{Original references}
\makeatletter
\newlabel{more-algebra-lemma-completion-regular}{{07NY}{source}{}{sourcepointer-6}{}}
\newlabel{more-algebra-lemma-completion-regular@cref}{{[originalreference][0][]07NY}{[1][0][]source}{}{}{}}
\makeatother
\crefname{originalreference}{original reference}{original references}
\Crefname{originalreference}{Original reference}{Original references}
\makeatletter
\newlabel{more-algebra-lemma-dominate-two-surjections}{{09Q8}{source}{}{sourcepointer-7}{}}
\newlabel{more-algebra-lemma-dominate-two-surjections@cref}{{[originalreference][0][]09Q8}{[1][0][]source}{}{}{}}
\makeatother
\crefname{originalreference}{original reference}{original references}
\Crefname{originalreference}{Original reference}{Original references}
\makeatletter
\newlabel{more-algebra-lemma-embed-map-Noetherian-complete-local-rings}{{07NN}{source}{}{sourcepointer-8}{}}
\newlabel{more-algebra-lemma-embed-map-Noetherian-complete-local-rings@cref}{{[originalreference][0][]07NN}{[1][0][]source}{}{}{}}
\makeatother
\crefname{originalreference}{original reference}{original references}
\Crefname{originalreference}{Original reference}{Original references}
\makeatletter
\newlabel{more-algebra-lemma-join-koszul-regular-sequences}{{0669}{source}{}{sourcepointer-9}{}}
\newlabel{more-algebra-lemma-join-koszul-regular-sequences@cref}{{[originalreference][0][]0669}{[1][0][]source}{}{}{}}
\makeatother
\crefname{originalreference}{original reference}{original references}
\Crefname{originalreference}{Original reference}{Original references}
\makeatletter
\newlabel{more-algebra-lemma-noetherian-finite-all-equivalent}{{09CC}{source}{}{sourcepointer-10}{}}
\newlabel{more-algebra-lemma-noetherian-finite-all-equivalent@cref}{{[originalreference][0][]09CC}{[1][0][]source}{}{}{}}
\makeatother
\title{The good-factorization kernel condition is independent of the factorization}
\author{The Stacks Project Authors}
\date{stacks-project:dpa.tex:local-snapshot-20261012; GNU FDL 1.2 or later}
\begin{document}
\maketitle
\noindent\textit{Editorial note.} No omission markers in-span; proof complete as written.\par\medskip
\noindent\textit{Allowed standard prerequisites (editorial).}\begin{itemize}
\item Domination of complete surjections. more-algebra-lemma-dominate-two-surjections, more-algebra-lemma-embed-map-Noetherian-complete-local-rings (cross-chapter, named). Produce the common dominating factorization S'' -> S x\_B S' used to reduce to comparable factorizations.
\item Regular sequence lifting. algebra-lemma-grothendieck-regular-sequence, more-algebra-lemma-join-koszul-regular-sequences, more-algebra-lemma-noetherian-finite-all-equivalent (cross-chapter, named). Drive the comparable-factorization step: lifting regular sequences through flat maps and joining them.
\end{itemize}\medskip
\ifdefined\hypertarget\hypertarget{sourcepointer-1}{}\fi
\noindent\textit{Source reference.} Original reference 00MB: focused-seeds/geometry/S9-dpa-ci-map-well-defined/ref-tags.txt, source lines 5--5. Navigation pointer only: the counted proof cites this statement by \textbackslash{}ref as a standing input; the deduction the counted proof has to make is carried out inside the counted span, and no part of that deduction is delegated to this pointer.\par\medskip
\ifdefined\hypertarget\hypertarget{sourcepointer-2}{}\fi
\noindent\textit{Source reference.} Original reference 00MA: focused-seeds/geometry/S9-dpa-ci-map-well-defined/ref-tags.txt, source lines 6--6. Navigation pointer only: the counted proof cites this statement by \textbackslash{}ref as a standing input; the deduction the counted proof has to make is carried out inside the counted span, and no part of that deduction is delegated to this pointer.\par\medskip
\ifdefined\hypertarget\hypertarget{sourcepointer-3}{}\fi
\noindent\textit{Source reference.} Original reference 00LM: focused-seeds/geometry/S9-dpa-ci-map-well-defined/ref-tags.txt, source lines 7--7. Navigation pointer only: the counted proof cites this statement by \textbackslash{}ref as a standing input; the deduction the counted proof has to make is carried out inside the counted span, and no part of that deduction is delegated to this pointer.\par\medskip
\ifdefined\hypertarget\hypertarget{sourcepointer-4}{}\fi
\noindent\textit{Source reference.} Original reference 00MG: focused-seeds/geometry/S9-dpa-ci-map-well-defined/ref-tags.txt, source lines 8--8. Navigation pointer only: the counted proof cites this statement by \textbackslash{}ref as a standing input; the deduction the counted proof has to make is carried out inside the counted span, and no part of that deduction is delegated to this pointer.\par\medskip
\ifdefined\hypertarget\hypertarget{sourcepointer-5}{}\fi
\noindent\textit{Source reference.} Original reference 00NR: focused-seeds/geometry/S9-dpa-ci-map-well-defined/ref-tags.txt, source lines 9--9. Navigation pointer only: the counted proof cites this statement by \textbackslash{}ref as a standing input; the deduction the counted proof has to make is carried out inside the counted span, and no part of that deduction is delegated to this pointer.\par\medskip
\ifdefined\hypertarget\hypertarget{sourcepointer-6}{}\fi
\noindent\textit{Source reference.} Original reference 07NY: focused-seeds/geometry/S9-dpa-ci-map-well-defined/ref-tags.txt, source lines 10--10. Navigation pointer only: the counted proof cites this statement by \textbackslash{}ref as a standing input; the deduction the counted proof has to make is carried out inside the counted span, and no part of that deduction is delegated to this pointer.\par\medskip
\ifdefined\hypertarget\hypertarget{sourcepointer-7}{}\fi
\noindent\textit{Source reference.} Original reference 09Q8: focused-seeds/geometry/S9-dpa-ci-map-well-defined/ref-tags.txt, source lines 11--11. Navigation pointer only: the counted proof cites this statement by \textbackslash{}ref as a standing input; the deduction the counted proof has to make is carried out inside the counted span, and no part of that deduction is delegated to this pointer.\par\medskip
\ifdefined\hypertarget\hypertarget{sourcepointer-8}{}\fi
\noindent\textit{Source reference.} Original reference 07NN: focused-seeds/geometry/S9-dpa-ci-map-well-defined/ref-tags.txt, source lines 12--12. Navigation pointer only: the counted proof cites this statement by \textbackslash{}ref as a standing input; the deduction the counted proof has to make is carried out inside the counted span, and no part of that deduction is delegated to this pointer.\par\medskip
\ifdefined\hypertarget\hypertarget{sourcepointer-9}{}\fi
\noindent\textit{Source reference.} Original reference 0669: focused-seeds/geometry/S9-dpa-ci-map-well-defined/ref-tags.txt, source lines 13--13. Navigation pointer only: the counted proof cites this statement by \textbackslash{}ref as a standing input; the deduction the counted proof has to make is carried out inside the counted span, and no part of that deduction is delegated to this pointer.\par\medskip
\ifdefined\hypertarget\hypertarget{sourcepointer-10}{}\fi
\noindent\textit{Source reference.} Original reference 09CC: focused-seeds/geometry/S9-dpa-ci-map-well-defined/ref-tags.txt, source lines 14--14. Navigation pointer only: the counted proof cites this statement by \textbackslash{}ref as a standing input; the deduction the counted proof has to make is carried out inside the counted span, and no part of that deduction is delegated to this pointer.\par\medskip
\begin{lemma}
\label{lemma-ci-map-well-defined}
Let $A \to B$ be a local homomorphism of Noetherian complete local rings.
The following are equivalent
\begin{enumerate}
\item for some good factorization $A \to S \to B$ the kernel of
$S \to B$ is generated by a regular sequence, and
\item for every good factorization $A \to S \to B$ the kernel of
$S \to B$ is generated by a regular sequence.
\end{enumerate}
\end{lemma}

\begin{proof}
Let $A \to S \to B$ be a good factorization.
As $B$ is complete we obtain a factorization
$A \to S^\wedge \to B$ where $S^\wedge$ is the completion of $S$.
Note that this is also a good factorization:
The ring map $S \to S^\wedge$ is flat
(Algebra, Lemma \ref{algebra-lemma-completion-flat}),
hence $A \to S^\wedge$ is flat.
The ring $S^\wedge/\mathfrak m_A S^\wedge = (S/\mathfrak m_A S)^\wedge$
is regular since $S/\mathfrak m_A S$ is regular
(More on Algebra, Lemma \ref{more-algebra-lemma-completion-regular}).
Let $f_1, \ldots, f_r$ be a minimal sequence of generators
of $\Ker(S \to B)$. We will use without further mention
that an ideal in a Noetherian local ring is generated by a regular
sequence if and only if any minimal set of generators is a
regular sequence. Observe that $f_1, \ldots, f_r$
is a regular sequence in $S$ if and only if $f_1, \ldots, f_r$
is a regular sequence in the completion $S^\wedge$ by
Algebra, Lemma \ref{algebra-lemma-flat-increases-depth}.
Moreover, we have
$$
S^\wedge/(f_1, \ldots, f_r)R^\wedge =
(S/(f_1, \ldots, f_n))^\wedge = B^\wedge = B
$$
because $B$ is $\mathfrak m_B$-adically complete (first equality by
Algebra, Lemma \ref{algebra-lemma-completion-tensor}).
Thus the kernel of $S \to B$ is generated by a regular sequence
if and only if the kernel of $S^\wedge \to B$ is generated by a
regular sequence.
Hence it suffices to consider good factorizations where $S$ is complete.

\medskip\noindent
Assume we have two factorizations $A \to S \to B$ and
$A \to S' \to B$ with $S$ and $S'$ complete. By
More on Algebra, Lemma \ref{more-algebra-lemma-dominate-two-surjections}
the ring $S \times_B S'$ is a Noetherian complete local ring.
Hence, using More on Algebra, Lemma
\ref{more-algebra-lemma-embed-map-Noetherian-complete-local-rings}
we can choose a good factorization $A \to S'' \to S \times_B S'$
with $S''$ complete. Thus it suffices to show:
If $A \to S' \to S \to B$ are comparable good factorizations,
then $\Ker(S \to B)$ is generated by a regular sequence
if and only if $\Ker(S' \to B)$ is generated by a regular sequence.

\medskip\noindent
Let $A \to S' \to S \to B$ be comparable good factorizations.
First, since $S'/\mathfrak m_R S' \to S/\mathfrak m_R S$ is
a surjection of regular local rings, the kernel is generated
by a regular sequence
$\overline{x}_1, \ldots, \overline{x}_c \in
\mathfrak m_{S'}/\mathfrak m_R S'$
which can be extended to a regular system of parameters for
the regular local ring $S'/\mathfrak m_R S'$, see
(Algebra, Lemma \ref{algebra-lemma-regular-quotient-regular}).
Set $I = \Ker(S' \to S)$. By flatness of $S$ over $R$ we have
$$
I/\mathfrak m_R I =
\Ker(S'/\mathfrak m_R S' \to S/\mathfrak m_R S) =
(\overline{x}_1, \ldots, \overline{x}_c).
$$
Choose lifts $x_1, \ldots, x_c \in I$. These lifts form a regular sequence
generating $I$ as $S'$ is flat over $R$, see
Algebra, Lemma \ref{algebra-lemma-grothendieck-regular-sequence}.

\medskip\noindent
We conclude that if also $\Ker(S \to B)$ is generated by a
regular sequence, then so is $\Ker(S' \to B)$, see
More on Algebra, Lemmas
\ref{more-algebra-lemma-join-koszul-regular-sequences} and
\ref{more-algebra-lemma-noetherian-finite-all-equivalent}.

\medskip\noindent
Conversely, assume that $J = \Ker(S' \to B)$ is generated
by a regular sequence. Because the generators $x_1, \ldots, x_c$
of $I$ map to linearly independent elements of
$\mathfrak m_{S'}/\mathfrak m_{S'}^2$ we see that
$I/\mathfrak m_{S'}I \to J/\mathfrak m_{S'}J$ is injective.
Hence there exists a minimal system of generators
$x_1, \ldots, x_c, y_1, \ldots, y_d$ for $J$.
Then $x_1, \ldots, x_c, y_1, \ldots, y_d$ is a regular sequence
and it follows that the images of $y_1, \ldots, y_d$ in $S$
form a regular sequence generating $\Ker(S \to B)$.
This finishes the proof of the lemma.
\end{proof}

\end{document}
